%  -----------------------------------------------------------------------------
%  Author         : Bimalka Piyaruwan Thalagala
%  GitHub         : https://github.com/bimalka98
%  Date Created   : 01.09.2020
%  Last Modified  : 15.02.2020
%  -----------------------------------------------------------------------------

\documentclass[a4paper,11pt]{article}%,twocolumn
%% packages


\usepackage{amsmath} % needed for command eqref
\usepackage{amssymb} % needed for math fonts
\usepackage[colorlinks=true,breaklinks]{hyperref} % needed for creating hyperlinks in the document, the option colorlinks=true gets rid of the awful boxes, breaklinks breaks lonkg links (list of figures), and ngerman sets everything for german as default hyperlinks language
\usepackage[hyphenbreaks]{breakurl} % ben�tigt f�r das Brechen von URLs in Literaturreferenzen, hyphenbreaks auch bei links, die �ber eine Seite gehen (mit hyphenation).
\usepackage{xcolor}
\definecolor{c1}{rgb}{0,0,1} % blue
\definecolor{c2}{rgb}{0,0.3,0.9} % light blue
\definecolor{c3}{rgb}{0.3,0,0.9} % red blue
\hypersetup{
    linkcolor={c1}, % internal links
    citecolor={c2}, % citations
    urlcolor={c3} % external links/urls
}
%\usepackage{cite} % needed for cite
\usepackage[square,authoryear]{natbib} % needed for cite and abbrvnat bibliography style
\usepackage[nottoc]{tocbibind} % needed for displaying bibliography and other in the table of contents
\usepackage{graphicx} % needed for \includegraphics 
\usepackage{longtable} % needed for long tables over pages
\usepackage{bigstrut} % needed for the command \bigstrut
\usepackage{enumerate} % needed for some options in enumerate
%\usepackage{todonotes} % needed for todos
\usepackage{makeidx} % needed for creating an index
\makeindex
\usepackage{gensymb}
\usepackage{url}
\usepackage{psfrag}
\usepackage{multirow}
\usepackage{subfigure}
\usepackage{epstopdf}

%% page settings

\usepackage[top=20mm, bottom=20mm,left=15mm,right=15mm]{geometry} % needed for page border settings
\parindent=0mm % for space of first line of new text block
\sloppy % for writing with hyphenless justification (tries to)
\hyphenation{} % use hyphenation of tolerance parametershttp://www.jr-x.de/publikationen/latex/tipps/zeilenumbruch.html
\hyphenpenalty=10000
\exhyphenpenalty=10000
\usepackage{fancyhdr} % needed for head and foot options
%% my macros

%% Text fomats
\newcommand{\tbi}[1]{\textbf{\textit{#1}}}

%% Math fonts
\newcommand{\bbA}{\mathbb{A}}
\newcommand{\bbB}{\mathbb{B}}
\newcommand{\bbC}{\mathbb{C}}
\newcommand{\bbD}{\mathbb{D}}
\newcommand{\bbE}{\mathbb{E}}
\newcommand{\bbF}{\mathbb{F}}
\newcommand{\bbG}{\mathbb{G}}
\newcommand{\bbH}{\mathbb{H}}
\newcommand{\bbI}{\mathbb{I}}
\newcommand{\bbJ}{\mathbb{J}}
\newcommand{\bbK}{\mathbb{K}}
\newcommand{\bbL}{\mathbb{L}}
\newcommand{\bbM}{\mathbb{M}}
\newcommand{\bbN}{\mathbb{N}}
\newcommand{\bbO}{\mathbb{O}}
\newcommand{\bbP}{\mathbb{P}}
\newcommand{\bbQ}{\mathbb{Q}}
\newcommand{\bbR}{\mathbb{R}}
\newcommand{\bbS}{\mathbb{S}}
\newcommand{\bbT}{\mathbb{T}}
\newcommand{\bbU}{\mathbb{U}}
\newcommand{\bbV}{\mathbb{V}}
\newcommand{\bbW}{\mathbb{W}}
\newcommand{\bbX}{\mathbb{X}}
\newcommand{\bbY}{\mathbb{Y}}
\newcommand{\bbZ}{\mathbb{Z}}
\usepackage[ framed, numbered]{matlab-prettifier}%framed,%
\usepackage{listings}

\begin{document}

\begin{titlepage}
\center % Center everything on the page

%-------------------------------------------------------------------------------------
%	HEADING SECTIONS
%------------------------------------------------------------------------------------
\textbf{\large Department of Electronic and Telecommunication Engineering}\\[0.5cm]
\textbf{\Large University of Moratuwa, Sri Lanka}\\[1cm]
\textbf{\large EN2570 - Digital Signal Processing}\\[2cm]
\includegraphics[width=0.3\textwidth]{figures/uomlogo}\\[2cm]

	
%-------------------------------------------------------------------------------------
%	TITLE SECTION
%------------------------------------------------------------------------------------
\textbf{\Huge Design of an FIR Digital Filter for Prescribed Specifications}\\[0.2cm]
{\Large (Using the windowing method in conjunction with the Kaiser window)}\\[0.5cm]
\textbf{\Large Project Report}\\[5cm]


%----------------------------------------------------------------------------------------
%	MEMBERS SECTION
%----------------------------------------------------------------------------------------

\vfill
\textbf{\large Submitted by}\\[0.5cm]
\begin{minipage}{0.2\textwidth}
	\begin{flushleft}
		{\large Thalagala B.P.}\\[4mm]		
	\end{flushleft}
\end{minipage}
\hspace{2mm}
\begin{minipage}{0.2\textwidth}
	\begin{flushright}
		{\large 180631J }\\[4mm]
	\end{flushright}
\end{minipage}\\[1.5cm]

%----------------------------------------------------------------------------------------
%	DATE SECTION
%----------------------------------------------------------------------------------------

\textbf{\large Submitted on}\\[0.5cm]
\textbf{\Large \today} % Date, change the \today to a set date if you want to be precise

%----------------------------------------------------------------------------------------

\vfill % Fill the rest of the page with whitespace

\end{titlepage}
\tableofcontents

\begin{center}
	\textbf{\textit{* PDF is clickable}}
\end{center}



\textit{Note:}\\
\textit{Additionally all the materials related to Task can also be found at \url{}}



\pagebreak
%%-----------------------------------------------------------------------
\section{Modeling the RF propagation Using Matlab}
%%-----------------------------------------------------------------------
\subsection{Relationship between Free Space Path Loss and Frequency}

\textit{Consider  following meanings for the parameters}\\

\begin{tabular}{l l }
	$P_{RX}$ & = Received Power at the Receiving Antenna\\
	$P_{TX}$ & = Transmitted Power at the Transmitting Antenna\\
	$f$ & = Frequency of the wave in Hz\\
	$f_{GHz}$ & = Frequency of the wave in GHz\\
	$d$& = Distance between the antennas in m\\
	$d_{km}$& = Distance between the antennas in km\\
	$G_{TX}$& = Directive gain of the Transmitter\\
	$G_{RX}$& = Directive gain of the Receiver\\
	$c$& = Velocity of the electromagnetic waves in a vacuum\\

\end{tabular}\\[1cm]


\[
\begin{split}
10.\log_{10}(L) &= 10.\log_{10}(\frac{(4\pi.f.d)^2}{c^2})\\
L_{dB}& = 10.\log_{10}((4\pi.f.d)^2) - 10.\log_{10}(c^2)\\
&=20.\log_{10}(4\pi.f.d)-20.\log_{10}(c)\\
&=20.\log_{10}(4\pi)-20.\log_{10}(c) + 20.\log_{10}(f) + 20.\log_{10}(d)\\
&=20.\log_{10}(\frac{4\pi}{c}) + 20.\log_{10}(f) + 20.\log_{10}(d)\\
&= -147.5522168 + 20.\log_{10}(f_{GHz}.10^9) + 20.\log_{10}(d_{km}.10^3)\\
& = -147.5522168 + 20.\log_{10}(10^9)+ 20.\log_{10}(f_{GHz}) + 20.\log_{10}(10^3) + 20.\log_{10}(d_{km})\\
&= -147.5522168 + 180+ 20.\log_{10}(f_{GHz}) + 60 + 20.\log_{10}(d_{km})\\
&= -147.5522168 + 240+ 20.\log_{10}(f_{GHz}) + 20.\log_{10}(d_{km})\\
&= +92.44778322+20.\log_{10}(f_{GHz}) + 20.\log_{10}(d_{km})
\end{split}
\]

\textbf{\textit{Note : Axes of the following plots are given in the logarithmic scale and range of frequency was chosen from 50 GHz to 1000 GHz since some of the ITU-R models are only defined in the 10 GHz-1000 GHz range.}}

%\begin{figure}[!h]
%	\centering
%	\includegraphics[scale=0.35]{figures/FreeSpacePL.png}
%	\caption{Relationship between Free Space Path Loss and Frequency}
%\end{figure}



%%-----------------------------------------------------------------------
\pagebreak
\subsection{Rain attenuation, Fog attenuation and Atmospheric gas attenuation with Frequency}

\textit{Note : For the generation of following plots three of the Matlab built-in functions, namely \textbf{rainpl()\cite{matlab}, gaspl()\cite{matlab}, fogpl()\cite{matlab}} which are developed according to the ITU-R P Series recommendations were used and links for their documentations are given at the Reference section.}


\subsection{Variation of the Signal Power with the Distance}

\begin{tabular}{l r}
Chosen Carrier frequency & 50 GHz\\
Transmission power & 50 kW or 47 dB\\
Cable loss at Transmitter & 3 dB\\
Transmitter Gain & 30 dB\\
Receiver Gain &24.77 dB\\
Cable loss at Receiver &4 dB\\
Total Path Loss & Varies with Distance\\
\end{tabular}\\[1cm]

\subsection{Codes}


\lstinputlisting[basicstyle = \mlttfamily\scriptsize , style = Matlab-editor]{code/RFPL.m}


\pagebreak
\bibliographystyle{plain}
\bibliography{refer}

%---------------------------------------------------------------------------
\end{document}
